\section{From a successful answer to an enduring institution}
\label{sec:institution}
The economically important question is not only whether an intelligence can produce an excellent result. It is whether an owner can reliably obtain further valuable results after the underlying model, operating environment, commercial relationship, or mission requirement changes. A system that wins once but cannot be renewed may be less useful than a slightly weaker system that can preserve, adapt, and re-establish what made its work valuable.

This paper develops \emph{SUCCESSOR $\Omega$} as a proposed institution of mission intelligence. The institution is not a particular model. It is the maintained arrangement through which legitimate knowledge, executable methods, evaluation, and bounded authority serve an enduring purpose. Models are replaceable contributors. They are not the institution's identity, its accumulated evidence, or its right to act.

The founding distinction comes from Transformational Intelligence: the authorized purpose, the current formulation of a problem, and the instrument used to evaluate a result are different objects \citep{boucherTI}. A method can change without changing the purpose. A new evaluator requires separate justification. A better-looking score cannot authorize the builder to redefine success. Here we apply that distinction across time: the institution should preserve what remains useful without preserving every assumption, dependency, or permission of the system that originally produced it.

\begin{insight}[The continuity principle]
Continuity is not the refusal to change. It is the ability to change components without losing the mission, legitimate knowledge, or accountability for what has been done.
\end{insight}

\subsection{What this paper adds}
The earlier FLYWHEEL $\Omega$ studies how production and verification change feasible economic interactions, including the cash timing that finances subsequent work \citep{boucherFlywheel}. SUPER ALPHA connects viable productive arrangements to value capture and successive capability construction \citep{boucherSuper}. The present paper addresses a different bottleneck: \emph{what survives replacement, and when paying to preserve it is worthwhile}.

Mission-relative revalidation in the earlier Transformational Autonomous Enterprises separates reusable generic checks from fresh mission validation; its direct specialist remains a strong comparator \citep{boucherEnterprise}. Out-Reflect constructs fault-bounded self-correcting reference evidence within a declared mathematical family \citep{boucherReflect}. We inherit those distinctions rather than treating a retained certificate or agreeing models as sufficient truth.

We make four concrete additions. First, renewal-adjusted Alpha evaluates an advantage over the lifecycle in which it must remain usable, rather than at one model release. Second, a context-bound succession contract separates portable knowledge from current evidence and authority. Third, an exact finite study identifies when a portability investment beats repeated rebuilding, when it does not, and why a competent ordinary planner can reproduce the result. Fourth, a physical-expansion contract makes the conversion of cash and knowledge into commissioned capacity inspectable rather than automatic.

These are a framework, conditional mathematical results, and newly executed software evidence. The full institutional architecture and industrial programme remain proposed. The public conceptual vocabulary has been synthesized from the author's programme materials; the private working conversations and supplied visual source files are not redistributed. Earlier public papers provide the stable citable foundation. This is a research edition, not a claim of independent peer review.

\subsection{The opportunity is more than continuity software}
The consequential target is the construction of productive capabilities that organizations cannot currently sustain economically. A laboratory might discover a useful process but repeatedly lose the ability to qualify it after instruments change. An industrial operator might possess excellent models but no portable, testable representation of operating conditions. A group of specialists might solve individual tasks yet fail at the interfaces required for a complete customer outcome.

The institution should therefore search over representations, experimental instruments, interfaces, production sequences, commercial commitments, and coordination arrangements. It may discover that the best move is a conventional method, a smaller system, or no investment. An institution exists to make valuable work possible, not to make itself indispensable.

The long-horizon ambition is correspondingly larger than retaining a chat history. It is to turn successful intelligence into accumulated scientific and productive means, and to use genuinely earned resources to build what the next consequential objective requires. Every step must preserve the difference between a plausible opportunity and a commissioned capability.

\section{The architecture of ownership}
\label{sec:ownership}
Represent a mission institution at time $t$ by
\begin{equation}
\mathcal I_t=(J,\mathcal M, K_t,X_t,H_t,E_t,A_t,L_t),
\end{equation}
where $J$ is its persistent identity, $\mathcal M$ the owner-authorized mission charter, $K_t$ retained knowledge, $X_t$ executable methods and instruments, $H_t$ operating history, $E_t$ current scoped evidence, $A_t$ active authority, and $L_t$ resources and obligations. A model-provider configuration is a replaceable dependency of $X_t$, not the definition of $J$.

The distinction between $K_t$ and $E_t$ is essential. The knowledge that a method once passed an evaluation may remain true. The claim that it is qualified \emph{now}, for another environment and mission, may be false. Likewise, knowing how to perform an action does not grant permission to perform it.

\subsection{Ownership must have an operational test}
Customer ownership is a proposed product requirement, not a legal conclusion established by this manuscript. Its technical acceptance test has four parts: the customer can export the authorized institutional assets; restore them in a separately controlled environment; replace a model or supplier; and recover the relevant bounded behavior under fresh evaluation. Rights to data, software, model weights, and external services must be established separately. Calling a bundle portable does not make a proprietary dependency portable.

The exported assets should include versioned methods, supported input and output contracts, scoped evidence history, provenance, failure conditions, routing procedures, and the real cost of maintaining the bundle. Credentials and active grants are not portable knowledge. A receiving installation begins with no inherited operating authority.

An owner may decide not to preserve an asset whose maintenance exceeds its future value. The institution must support retirement as well as accumulation. Compact, decision-relevant knowledge can be preferable to a larger archive that is expensive, stale, or impossible to retrieve reliably. Historical losses and obligations cannot disappear when operational material is compacted.

\subsection{A capability capsule is a hypothesis about reuse}
A useful capsule contains an executable artifact, the semantics of its inputs and outputs, prerequisites, dependency identities, prior test evidence, known counterexamples, and a proposed successor use. It does not contain a universal permission to reuse the earlier conclusion. A capsule may be technically intact while inappropriate for its next task.

We distinguish four relationships between campaigns. A predecessor may \emph{finance} a successor, \emph{equip} it with a useful tool, \emph{enable} a construction that would otherwise be inaccessible within matched resources, or \emph{improve discovery} itself. These relationships are not interchangeable. A later result is not evidence of any of them merely because it occurs later.

The strongest institutional asset may be a constructor: a tested method for producing new instruments under changed conditions. But a narrative explaining a successful instrument is not automatically such a constructor. It must generate executable candidates and survive fresh comparisons, including ordinary reuse of the same substantive information.

\section{Alpha, Gym, Specialist ASI, SUCCESSOR \texorpdfstring{$\Omega$}{Omega}}
\label{sec:stages}
The four names describe distinct functions. They need not impose a rigid four-agent implementation, nor require every mission to traverse every stage again. A qualified conventional tool may already be the best route. What cannot be skipped is the evidence or authority required for the actual action.

\subsection{Alpha: find the advantage worth building}
Alpha begins with a beneficiary and a consequential outcome, not a favorite architecture. Identify the current complete method and the best credible alternative, including self-service, other autonomous providers, conventional software, cooperation, and non-entry where relevant. Investigate what prevents the desired outcome from being delivered economically, reliably, or at the necessary scale.

An Alpha hypothesis specifies the proposed change, its causal mechanism, the outcome that remains protected, its full lifecycle burden, and a decisive failure condition. A lower token bill is not enough if acquisition, verification, integration, delay, or maintenance overwhelms the saving. A technically ordinary component can be part of an extraordinary complete arrangement; novelty is not an exemption from comparison.

The research executive may search backward from the outcome toward missing means, or forward from available capabilities toward newly reachable outcomes. It may transfer a structure from another domain, construct a missing correspondence, or remove the dependency that made the correspondence seem necessary. The choice among those routes is itself an object of investigation, not a fixed strategy imposed on every task.

\subsection{Gym: develop the capability before trusting it}
A Gym is a mission development environment built from authorized evidence, tools, constraints, and experimental history. It supports proposal, construction, simulation, stress testing, comparison, and correction. Its output is a candidate and a development record, not a certificate of deployment readiness.

The Gym must represent failure modes that matter: missing observations, wrong initial hypotheses, unaffordable tests, incompatible interfaces, misleading references, absent demand, delayed collection, and the possibility that the incumbent is already sufficient. Difficulty should not be manufactured merely to make the new architecture look useful.

Simulation can identify a discriminating physical experiment. It cannot substitute for that experiment when the claim depends on unknown physical behavior. NASA's readiness distinctions between analytical, laboratory, relevant-environment, and operational demonstrations are a useful precedent for keeping those steps separate \citep{nasaTRL}. We adopt the distinction, not an automatic aerospace certification for this programme.

\subsection{Specialist ASI: an earned, bounded designation}
In this programme, \spec{} is a proposed mission-specific designation. It is not a standardized scientific category or a claim of general superintelligence. An eligible claim must identify the exact candidate, mission distribution, protected test protocol, comparators, resource accounting, practical improvement threshold, safety constraints, uncertainty, and validity conditions. The designation must never travel without that scope.

Freeze the candidate before protected confirmation. Compare it with the current approach and the best credible complete alternative under comparable permissions and resources. Evaluate accepted outcomes, total cost, time, robustness, and material risks. Multiple objectives require declared tradeoffs or non-inferiority conditions; a convenient weighted average must not conceal a disqualifying regression.

No experiment in this paper awards that designation to a model. The cash study evaluates supplied operating arrangements, not AI task performance. The local state machine tests acceptance logic, not the truth of an ASI claim. A label cannot supply the evidence it names.

\subsection{SUCCESSOR \texorpdfstring{$\Omega$}{Omega}: preserve, govern, renew}
Following sufficient evidence and explicit customer authorization, a capability may become part of the operating institution. The institution preserves legitimate learning, executes within a bounded grant, monitors drift, records consequential changes, and supports expiry, suspension, recovery, and renewal.

These functions form a feedback relationship, not an irreversible staircase. Operational evidence may return a capability to the Gym. A supplier replacement may require only a narrow justified requalification, or a complete rebuild. A changed mission may invalidate the original abstraction. A stopped capability can remain informative without remaining active.

\begin{figure}[htbp]
\centering
\begin{tikzpicture}[node distance=7mm,box/.style={draw=gold,fill=cream,rounded corners,align=center,text width=3.1cm,minimum height=17mm,font=\small},arr/.style={-{Stealth},thick,color=wine}]
\node[box] (a) {\textbf{Alpha}\\Worth constructing?};
\node[box,right=of a] (g) {\textbf{Gym}\\Candidate + history};
\node[box,right=of g] (s) {\textbf{Fresh evidence}\\Qualified in scope?};
\node[box,right=of s] (o) {\textbf{Owner grant}\\Bounded operation};
\draw[arr] (a)--(g);\draw[arr] (g)--(s);\draw[arr] (s)--(o);
\draw[arr] (o.south) -- ++(0,-.6) -| node[pos=.25,below,font=\small]{Renew, replace, retire, or investigate} (g.south);
\end{tikzpicture}
\caption{The four-stage vocabulary is implemented through distinct decisions. The development environment does not certify itself; evidence does not issue its own operating permission.}
\end{figure}
